\centerline{\bf \Huge Разностная последовательность}

\begin{flushright}
        \scriptsize{Эпизод 1. Нападение Ангела.}
\end{flushright}

\textit{Разностной последовательностью мы будем называть последовательность соседних разностей, а именно по последовательности $a_1, a_2, \ldots, a_n, \ldots$ можно построить новую последовательность $b_n=a_{n+1}-a_n$. После такой операции часть информации разумеется пропадет, но если запомнить, чему равняется $a_1$, то исходная последовательность однозначно восстанавливается. Такой трюк иногда позволяет переписать известное в нескольких других, более удобных терминах.}

\problem{1}
Конечная возрастающая последовательность $a_1, a_2, a_3, \ldots, a_n$ состоит из $n \geqslant 3$ различных натуральных чисел, причем при всех натуральных $k \leqslant n-2$ выполнено равенство $4 a_{k+2}=7 a_{k+1}-3 a_k$.

(а) Приведите пример такой последовательности для $n=5$. 

(b) Может ли в такой последовательности при некотором $n \geqslant 3$ выполняться равенство $a_n=4 a_2-3 a_1$?

(c) Какое наименьшее значение может принимать $a_1$, если $a_n=527$ ?

\problem{2}
Конечная последовательность $a_1, a_2, \ldots, a_n$ состоит из $n \geqslant 3$ не обязательно различных натуральных чисел, причем при всех натуральных $k \leqslant n-2$ выполнено равенство $a_{k+2}=2 a_{k+1}-a_k-2$. 

(a) Приведите пример такой последовательности при $n=6$, в которой $a_6=6$. 

(b) Может ли в такой последовательности некоторое натуральное число встретиться три раза?

(c) При каком наибольшем $n$ такая последовательность может состоять только из четных двузначных чисел?

\problem{3}
Последовательность $a_1, a_2, \ldots, a_n(n \geqslant 3)$ состоит из натуральных чисел, причем каждый член последовательности (кроме первого и последнего) больше среднего арифметического соседних (стоящих рядом с ним) членов. 

(а) Приведите пример такой последовательности, состоящей из четырех членов, сумма которых равна 50 . 

(b) Может ли такая последовательность состоять из шести членов и содержать два одинаковые числа?

(c) Какое наименьшее значение может принимать сумма членов такой последовательности при $n=10$ ?

\problem{4}
Последовательность положительных вещественных чисел $x_n$ такова, что $x_n^2<x_n-x_{n+1}$. Докажите, что $x_n<\frac{1}{n}$ для всех натуральных $n$.

\problem{5}
Петя хочет выписать все возможные последовательности из 100 натуральных чисел, в каждой из которых хотя бы раз встречается тройка, а любые два соседних члена различаются не больше, чем на 1 . Сколько последовательностей ему придётся выписать?

\problem{6}
Пусть $a_0, a_1, a_2, \ldots$ бесконечная последовательность действительных чисел, удовлетворяющих уравнению $a_n=\left|a_{n+1}-a_{n+2}\right|$ для всех $n \geqslant 0$, где $a_0$ и $a_1$ два различных положительных действительных числа. Может ли последовательность $a_0, a_1, a_2, \ldots$ быть ограниченной?